% ID: SEL-D0056-OA-E04
% ESTADO: cotejada. Segunda pasada visual y borrador compilado revisados.
% Fuente: Selectividad/sel_2011_matematicas_aplicadas/123-EXAMEN 1.pdf
% SHA-256: 69f8b84ed22f2b49074eafb95dbdce8a98114a4a2f02989b92f3f7add9355155
% Página: 1; opción A; ejercicio original 4.

El director de un banco afirma que la cantidad media de dinero extraído, por cliente, de un cajero automático de su sucursal no supera los 120 euros. Para contrastar esta hipótesis elige al azar 100 extracciones de este cajero y obtiene una media muestral de 130 euros. Se sabe que la cantidad de dinero extraído por un cliente en un cajero automático se distribuye según una ley Normal de media desconocida y desviación típica 67 euros.
\begin{enumerate}[label=\alph*)]
\item \textbf{(0.5 puntos)} Plantee el contraste de hipótesis asociado al enunciado.
\item \textbf{(1 punto)} Determine la región de aceptación, para un nivel de significación $\alpha=0.05$.
\item \textbf{(1 punto)} Con los datos muestrales tomados, \textquestiondown{}existe evidencia estadística para rechazar la hipótesis de este director, con el mismo nivel de significación anterior?
\end{enumerate}
